\documentclass[12pt]{article}

\usepackage[a4paper,margin=14mm,top=12mm,bottom=12mm]{geometry}
\usepackage[T1]{fontenc}
\usepackage[scaled=0.95]{helvet}
\usepackage{amsmath}
\usepackage{xcolor}
\usepackage{enumitem}

\pagestyle{empty}
\renewcommand{\familydefault}{\sfdefault}
\setlength{\parindent}{0pt}
\setlength{\parskip}{1.2mm}

\begin{document}

\begin{center}
  {\LARGE\bfseries Worked solutions: Savings accounts}\\[1.5mm]
  {\small All amounts are rounded to the nearest penny.}
\end{center}

\section*{Part 1 worked solutions}

For simple interest use $P(1+rt)$.  
For compound interest use $P(1+r)^t$.  
In part 1, every account adds interest once a year.

\begin{enumerate}[label=\textbf{\arabic*.}, leftmargin=4mm, itemsep=6mm, topsep=2mm]

\item \textbf{Amara has \pounds100 and saves for 1 year.}

Virgo cannot be used because it needs \pounds500.  
Tesbury cannot be used because it needs \pounds200.  
Metroline's \pounds300 maximum is fine for \pounds100.

\[
\text{Metroline: } 100(1+0.10)=110
\]
\[
\text{Coventree: } 100(1.05)=105,\qquad
\text{Starlight: } 100(1.05)=105,\qquad
\text{Chace: } 100(1.04)=104
\]

\textbf{Best account: Metroline Bank.}  
\textbf{Amara finishes with \pounds110.}

\item \textbf{Ben has \pounds200 and saves for 2 years.}

Metroline and Tesbury can both be used.  
Tesbury requires the money to stay in for at least 2 years, and Ben's plan is exactly 2 years.

\[
\text{Metroline: } 200+2(0.10\times 200)=200+40=240
\]
\[
\text{Tesbury: } 200+2(0.06\times 200)=200+24=224
\]
\[
\text{Starlight: } 200(1.05)^2=200(1.1025)=220.50
\]
\[
\text{Coventree: } 200+2(0.05\times 200)=220
\]
\[
\text{Chace: } 200(1.04)^2=200(1.0816)=216.32
\]

\textbf{Best account: Metroline Bank.}  
\textbf{Ben finishes with \pounds240.}

\item \textbf{Chloe has \pounds400 and saves for 4 years.}

Metroline cannot be used because \pounds400 is above its \pounds300 maximum.  
Virgo cannot be used because it needs \pounds500.  
Tesbury can be used because Chloe's money stays in for at least 2 years.

\[
\text{Tesbury: } 400+4(0.06\times 400)=400+96=496
\]
\[
\text{Starlight: } 400(1.05)^4=400(1.21550625)=486.20
\]
\[
\text{Coventree: } 400+4(0.05\times 400)=400+80=480
\]
\[
\text{Chace: } 400(1.04)^4=400(1.16985856)=467.94
\]

\textbf{Best account: Tesbury Bank.}  
\textbf{Chloe finishes with \pounds496.}

\item \textbf{Dev has \pounds1000 and saves for 3 years.}

Metroline cannot be used because \pounds1000 is above its \pounds300 maximum.  
Virgo can be used because the minimum is \pounds500.

\[
\text{Virgo: } 1000(1.08)^3=1000(1.259712)=1259.71
\]
\[
\text{Tesbury: } 1000+3(0.06\times 1000)=1000+180=1180
\]
\[
\text{Starlight: } 1000(1.05)^3=1000(1.157625)=1157.63
\]
\[
\text{Coventree: } 1000+3(0.05\times 1000)=1150
\]
\[
\text{Chace: } 1000(1.04)^3=1000(1.124864)=1124.86
\]

\textbf{Best account: Virgo Money.}  
\textbf{Dev finishes with \pounds1259.71.}

\item \textbf{Erin has \pounds500 and saves for 5 years.}

Metroline cannot be used because \pounds500 is above its \pounds300 maximum.  
Virgo can be used because \pounds500 meets its minimum.

\[
\text{Virgo: } 500(1.08)^5=500(1.4693280768)=734.66
\]
\[
\text{Tesbury: } 500+5(0.06\times 500)=500+150=650
\]
\[
\text{Starlight: } 500(1.05)^5=500(1.2762815625)=638.14
\]
\[
\text{Coventree: } 500+5(0.05\times 500)=625
\]
\[
\text{Chace: } 500(1.04)^5=500(1.2166529024)=608.33
\]

\textbf{Best account: Virgo Money.}  
\textbf{Erin finishes with \pounds734.66.}

\item \textbf{Fay has \pounds250, saves for 3 years, and is comparing Coventree and Starlight only.}

Coventree pays simple interest:

\[
250+3(0.05\times 250)=250+37.50=287.50
\]

Starlight pays compound interest:

\[
250(1.05)^3=250(1.157625)=289.40625\approx 289.41
\]

Difference:

\[
289.41-287.50=1.91
\]

\textbf{Fay should choose Starlight Bank.}  
\textbf{She is \pounds1.91 better off than with Coventree.}

\end{enumerate}

\section*{Part 2 worked solutions}

Remember:
\begin{itemize}[leftmargin=5mm, itemsep=1mm]
  \item Monza adds $0.5\%$ each month, so use $1.005$ per month.
  \item Halifield adds $2\%$ each half-year, so use $1.02$ per half-year.
  \item Nationworth charges \pounds2 per month, so charges \pounds24 per year.
  \item Trustbridge pays $2\%+4\%+6\%=12\%$ in total over 3 years, all on the original deposit.
\end{itemize}

\begin{enumerate}[label=\textbf{\arabic*.}, leftmargin=4mm, itemsep=6mm, topsep=2mm]

\item \textbf{Priya is 13, has \pounds200, and saves for 1 year.}

NatEast is available because she is aged 11--16 and \pounds200 is between \pounds50 and \pounds300.

Santandra requires at least 2 years, so it is not available.  
Lloydale does not allow withdrawals for 3 years, so it is not available.  
Nationworth needs a minimum of \pounds250, so it is not available.  
Trustbridge requires the full 3 years, so it is not available.

\[
\text{NatEast: } 200(1.09)=218
\]
\[
\text{Barcourt: } 200(1.05)=210
\]
\[
\text{Monza: } 200(1.005)^{12}=212.34
\]
\[
\text{Halifield: } 200(1.02)^2=200(1.0404)=208.08
\]

\textbf{Best account: NatEast Young Saver.}  
\textbf{Priya finishes with \pounds218.}

\item \textbf{Marcus is 30, has \pounds1000, and saves for 2 years.}

NatEast is not available because Marcus is too old.  
Lloydale and Trustbridge require 3 years, so they are not available.

\[
\text{Santandra: } 1000(1.08)^2=1000(1.1664)=1166.40
\]
\[
\text{Monza: } 1000(1.005)^{24}=1127.16
\]
\[
\text{Barcourt: } 1000(1.05)^2=1000(1.1025)=1102.50
\]
\[
\text{Halifield: } 1000(1.02)^4=1000(1.08243216)=1082.43
\]
\[
\text{Nationworth: } 1000(1.12)-48=1120-48=1072
\]

\textbf{Best account: Santandra Fixed Saver.}  
\textbf{Marcus finishes with \pounds1166.40.}

\item \textbf{Sofia is 28, has \pounds150, and saves for 3 years.}

Santandra is not available because it needs \pounds200.  
NatEast is not available because Sofia is too old.  
Nationworth is not available because it needs \pounds250.  

Lloydale can be used: the money stays in for 3 full years, so the \pounds25 bonus is paid.

\[
\text{Lloydale: } 150(1.03)^3+25
\]
\[
150(1.092727)=163.91,\qquad 163.91+25=188.91
\]

Compare with the other available accounts:

\[
\text{Monza: } 150(1.005)^{36}\approx 179.50
\]
\[
\text{Barcourt: } 150(1.05)^3=150(1.157625)=173.64
\]
\[
\text{Halifield: } 150(1.02)^6=150(1.126162419264)=168.92
\]
\[
\text{Trustbridge: } 150(1+0.02+0.04+0.06)=150(1.12)=168
\]

\textbf{Best account: Lloydale Bonus Builder.}  
\textbf{Sofia finishes with \pounds188.91.}

\item \textbf{Dan is 40, has \pounds2000, and saves for 2 years.}

NatEast is not available because Dan is too old.  
Lloydale and Trustbridge require 3 years, so they are not available.  
Monza cannot hold the full \pounds2000 because its maximum balance is \pounds1000.

\[
\text{Santandra: } 2000(1.08)^2=2000(1.1664)=2332.80
\]
\[
\text{Barcourt: } 2000(1.05)^2=2000(1.1025)=2205
\]
\[
\text{Nationworth: } 2000(1.12)-48=2240-48=2192
\]
\[
\text{Halifield: } 2000(1.02)^4=2000(1.08243216)=2164.86
\]

\textbf{Best account: Santandra Fixed Saver.}  
\textbf{Dan finishes with \pounds2332.80.}

\item \textbf{Aisha is 25, has \pounds500, and saves for 1 year.}

Santandra needs at least 2 years, so it is not available.  
NatEast is not available because Aisha is too old.  
Lloydale and Trustbridge require 3 years, so they are not available.

\[
\text{Monza: } 500(1.005)^{12}=500(1.0616778118645)=530.84
\]
\[
\text{Barcourt: } 500(1.05)=525
\]
\[
\text{Halifield: } 500(1.02)^2=500(1.0404)=520.20
\]
\[
\text{Nationworth: } 500(1.06)-24=530-24=506
\]

\textbf{Best account: Monza Monthly Saver.}  
\textbf{Aisha finishes with \pounds530.84.}

\item \textbf{Leo is 16, has \pounds300, and saves for 3 years.}

NatEast is available because Leo is aged 11--16 and \pounds300 is exactly its maximum.

\[
\text{NatEast: } 300(1+3\times0.09)=300(1.27)=381
\]
\[
\text{Santandra: } 300(1.08)^3=300(1.259712)=377.91
\]
\[
\text{Monza: } 300(1.005)^{36}\approx 359.00
\]
\[
\text{Lloydale: } 300(1.03)^3+25=300(1.092727)+25=352.82
\]
\[
\text{Barcourt: } 300(1.05)^3=300(1.157625)=347.29
\]
\[
\text{Halifield: } 300(1.02)^6=300(1.126162419264)=337.85
\]
\[
\text{Trustbridge: } 300(1+0.02+0.04+0.06)=300(1.12)=336
\]
\[
\text{Nationworth: } 300(1+3\times0.06)-36\times2=300(1.18)-72=354-72=282
\]

The two closest accounts are NatEast and Santandra:

\[
381-377.91=3.09
\]

\textbf{Best account: NatEast Young Saver with \pounds381.}  
\textbf{Santandra gives \pounds377.91, so the difference is \pounds3.09.}

\end{enumerate}

\section*{Think about}

\begin{itemize}[leftmargin=5mm, itemsep=3mm, topsep=1mm]

\item \textbf{Coventree versus Starlight.}

Both pay $5\%$. After one year they give exactly the same because there is no previous interest to compound. After that, Starlight's $5\%$ is applied to the growing balance, including interest already earned, while Coventree's $5\%$ is applied only to the original deposit. So Starlight always pays at least as much as Coventree, and pays more whenever the money is left in for more than one year.

\item \textbf{Is the biggest percentage always the best?}

No. In part 1, Metroline has the highest rate, $10\%$, but its maximum balance is only \pounds300. Chloe has \pounds400, Dev has \pounds1000, and Erin has \pounds500, so none of them can use Metroline. So the highest advertised percentage rate does not guarantee the best result; conditions such as minimum deposits, maximum balances, and time restrictions matter too.

\item \textbf{Why is Trustbridge not a good deal?}

Trustbridge pays $2\%$, then $4\%$, then $6\%$, but all are simple interest on the original deposit. Over 3 years the total interest is
\[
2\%+4\%+6\%=12\%,
\]
which is only $4\%$ per year on average. The money also has to stay in for the full 3 years, and there is no compounding. Several other accounts in part 2 beat it.

\item \textbf{Why does nobody choose Nationworth?}

Nationworth charges a \pounds2 fee every month, which is \pounds24 per year. On the minimum deposit of \pounds250, the annual simple interest is only
\[
0.06\times 250=\pounds15,
\]
so the fee is larger than the interest. Even on \pounds500, the gross interest is \pounds30 but the fees are \pounds24, leaving only \pounds6. You need over \pounds400 in the account before the $6\%$ interest even covers the annual fee.

\item \textbf{Is Halifield paying 2\% twice a year the same as 4\% once a year?}

No. Halifield adds interest twice a year, and the second $2\%$ is calculated on the balance after the first $2\%$ has been added. Over one year:
\[
1.02\times 1.02=1.0404.
\]
So $2\%$ twice a year is equivalent to $4.04\%$ once a year, slightly more than $4\%$.

\end{itemize}

\end{document}